Este trabajo no puede estar dedicado a nadie más que a mis compañeros del grado de matemáticas.

Estudiar por la UNED no es fácil, pues estamos solos ante el peligro (los libros ladrillo\footnote{Sí, me refiero a ti, \textit{Integración en variedades}, la que has liado tú solito.}). El único resquicio que nos queda a muchos es ayudarnos: compartir material, debatir, dudar, pensar, equivocarnos, rectificar, criticar, compartir éxitos y lamentos. Es gracias a los compañeros que crearon y mantienen los grupos de Telegram que muchas asignaturas complicadas se vuelven amenas por el mero hecho de poder discutir problemas con otros. También, los compañeros que han compartido de manera desinteresada sus materiales de estudio, sus opiniones, sus consejos, aportan su granito de arena para los que hemos venido tras ellos, y por lo tanto, recae en nosotros, ahora, aportar para los que vengan después. Es por ello que le dedico este trabajo a mis compañeros. A los que ya terminaron; a los que he tenido la oportunidad de conocer: los que me han ayudado, con los que he debatido, y a los que he ayudado; y a los estudiantes que están por venir. Sobre todo las últimas asignaturas de la carrera, cuando estas se vuelven más ariscas y hay menos actividad, la necesidad de tener compañeros a los que atenerse es fundamental. Es por ello que también estoy agradecido con los amigos del grado a los que el contacto ha logrado sobrevivir, y con los que espero seguir relacionándome.

\medskip

No puedo agradecer primero a nadie más que a mis padres, que me lo han dado todo y me siguen apoyando, a pesar de que muchas veces mis decisiones no sean las más acertadas; y a mi hermano, que siempre está ahí cuando se le necesita. En segundo lugar, a mis segundos padres, Fran y Ángela, mis profesores del grado superior de música y que me han cuidado y guiado cuando estaba perdido. En tercer lugar, a mi pareja Ana, quien, aunque le diga que intente leerse mi TFG, me argüirá de nuevo que no lo entiende. Ciertamente, se hace complicado exponer los resultados para alguien que no ha estudiado matemáticas. Para los matemáticos que estén leyendo esto, al menos darle un intento. Quizás no sea el mejor con el uso de las palabras, pero se hace lo que se puede. Por último, agradecer a mis amigos, que hacen la vida más amena.

\newpage

     \href{https://github.com/Jex-here/Trabajo-Final-de-Grado-Matematicas}{\textit{En este repositorio de github}} se encuentra el código látex del trabajo. Para descargarlo, si nunca usaste github, pulsa en CODE y DOWNLOAD ZIP, tal como se muestra en \ref{fig:github}. Luego, puedes subir el .zip a overleaf, o descomprimirlo y abrir la carpeta en tu entorno de \LaTeX \ de confianza.
\underline
    
    
\begin{figure}[H]
        \centering
        \includegraphics[width=1\linewidth]{imagenes/github.png}
        \caption{Explicación github}
        \label{fig:github}
\end{figure}
        

